\section{Hash Functions}\label{sec:2}

A hash function maps bit strings of an arbitrary finite length to
fixed-length hash values. On the one hand a hash function must be fast
to compute, but on the other it must also be collision-resistant, i.e.
it must be computationally infeasible to find a collision (that is a
pair of messages with the same hash value) \cite{ref11}.
Additionally, hash functions are intended to resist inversion (i.e.
finding a message from a given hash value).

The basic idea of cryptographic hash functions is that a hash value
serves as a compact representative image (sometimes called a
\emph{message digest}) of an input string, and can be used to uniquely
identify that string. There are several types of hash functions
\cite{ref6}\cite{ref25}. The best known are:

\begin{itemize}
\tightlist
\item
  Hash functions based on block ciphers;
\item
  Hash functions based on modular arithmetic;
\item
  Hash functions based on cellular automata;
\item
  Hash functions based on knapsack problem;
\item
  Hash functions based on algebraic matrices; and
\item
  Dedicated hash functions.
\end{itemize}

In this paper we focus on dedicated hash functions. Readers interested
in a non-mathematical treatment of hash functions and cryptography are
referred to \cite{ref28} and \cite{ref27}
respectively.

\subsection{Definitions}\label{sec:2.1}\label{definitions}

Perhaps the most widely known definitions for collision resistance,
preimage-resistance, and second-preimage resistance come from the
Handbook of Applied Cryptography (HAC) \cite{ref2} which is
based in part on the work of Preneel \cite{ref5} which is the
seminal work in this area, providing extensive treatment of
cryptographic hash functions according to the information theoretic,
complexity theoretic, and system based approaches. For our purposes we
will rely primarily upon his descriptions.

The following terms are basic to cryptography in general and hash
functions in particular.

\emph{Definition 1}: A \textbf{one-way hash function} (OWHF) is a
function \emph{h} satisfying the following conditions:

\begin{enumerate}
\def\labelenumi{\arabic{enumi}.}
\tightlist
\item
  The argument X can be of arbitrary length and the result h(X) has a
  fixed length of n bits.
\item
  Given h and X, the computation of h(X) must be ``easy''.
\item
  \emph{The hash function must be one-way in the sense that given a Y in
  the image of h, it is ``hard'' to find a message X such that h(X) = Y
  (preimage resistant) and given X and h(X) it is ``hard'' to find a
  message \(X'\ne X\) such that \(h(X')=h(X)\) (second-preimage
  resistant).}
\end{enumerate}

\emph{Definition 2}: A \textbf{collision resistant hash function} (CRHF)
is a one-way hash function with the additional property that the hash
function must be collision resistant: this means that it is ``hard'' to
find two distinct messages that hash to the same result. That is, given
any two distinct inputs \(X,X'\) it is computationally infeasible to
hash to the same output such that \(h(X')=h(X).\)

Note that ``preimage resistant'' is alternatively referred to as
\emph{one-way}; ``(second) preimage resistant'' is referred to as
\emph{weak collision resistant}; and ``collision resistant'' is referred
to as \emph{strong collision resistant}. Further mathematical rigor to
the definition of these terms can be found in \cite{ref2} and
\cite{ref17}.

\emph{Definition 3}: An \emph{n}-bit hash function has \textbf{ideal
security} if it satisfies the following conditions:

\begin{enumerate}
\def\labelenumi{\arabic{enumi}.}
\tightlist
\item
  \emph{Given a hash output, producing a preimage or second preimage
  requires \textasciitilde{} \(2^n\) operations.}
\item
  \emph{Producing a collision requires \textasciitilde{} \(2^{n/2}\)
  operations.}
\end{enumerate}

\emph{Definition 4}: \textbf{Manipulation detection code} (MDC) is the
umbrella term used to describe the class of hash functions that are
unkeyed.

The tradeoff that must be considered when choosing a specific algorithm
in this class is the tension between efficiency of computation versus
degree of security. The two are often inversely related. Because the
speed of a hash function is typically of critical importance, many block
ciphers -- while providing excellent security - are too slow to be
practical. This gave rise to specialized hash functions, which are
designed from scratch with speed as a principal consideration.

\emph{Definition 5}: \textbf{Message authentication code} (MAC) is the
umbrella term used to describe the class of hash functions that are
keyed. While the introduction of a key adds state, it has the positive
effect of allowing for expanded use of the hash function, with
applications in areas such as digital signature authentication. For
example, Centera uses a keyed MAC to validate operations between
management station(s) and the Centera server.

A MAC is a function which takes the secret key \emph{k} (shared between
Alice and Bob) and the message \emph{m} to return a tag
MAC\textsubscript{\emph{k}}(\emph{m}). Here, Eve represents the
adversary (eavesdropper). Eve sees a sequence
(\emph{m}\textsubscript{\emph{1}}\emph{,a}\textsubscript{\emph{1}}),
(\emph{m}\textsubscript{\emph{2}}\emph{,a}\textsubscript{\emph{2}}),\ldots,(\emph{m}\textsubscript{\emph{q}}\emph{,a}\textsubscript{\emph{q}})
of pairs of messages and their corresponding tags (i.e.,
\emph{a}\textsubscript{\emph{i~}}\emph{=}
MAC\textsubscript{\emph{k}}\emph{(m}\textsubscript{\emph{i}}))
transmitted between the parties. Eve breaks the MAC if she can find a
message \emph{m}, not included among
\emph{m}\textsubscript{\emph{1}}\emph{,m}\textsubscript{\emph{2}}\emph{,,\ldots,m}\textsubscript{\emph{q}},
together with its corresponding valid authentication tag \emph{a} =
MAC\textsubscript{\emph{k}}(\emph{m}). Eve's success probability is the
probability that she breaks the MAC.

Hash functions are generally not keyed. However, there is interest in
using cryptographically strong hash functions as the basis for MAC
schemes. Bellare et al.~\cite{ref16} describe a means of using
an iterative hash function as the basis for a MAC by keying the hash
function via the initial vector, instead of the more common mechanism of
inputting the key as part of the data hashed by the function.

\subsection{Iterated Hash Functions}\label{sec:2.2}\label{iterated-hash-functions}

Most known hash functions (including MD5 and SHA) are based on a
compression function with fixed size input; they process every message
block in a similar way. The information is divided into \emph{t} blocks
\emph{X}\textsubscript{\emph{1}} through
\emph{X}\textsubscript{\emph{t}}. If the total number of bits is not a
multiple of the block length, the information is padded to the required
length (using a so-called \emph{padding rule}). The hash function can
then be described as follows:

\begin{equation}\label{eq:2.1}
\begin{aligned}H_0&=IV\\H_i&=f(X_i,H_{i-1})\quad i=1,2,\ldots,t\\h(X)&=g(H_t)\end{aligned}
\end{equation}

Where \emph{IV} is the abbreviation for Initial Value (or Initial
Vector); \(H_{i-1}\) serves as the \emph{n}-bit chaining variable
between stage \emph{i} -- 1 and stage \emph{i}; and the result of the
hash function is denoted with \emph{h(X).} The function \emph{f} is
called the \emph{round} or \emph{compression} function, and the function
\emph{g} is called the \emph{output transformation}. It is often omitted
(that is, \emph{g} is often the identity function); in the case of
Centera, \emph{g} maps bit strings to a base32-notation (digits 0-9 and
A-V), which has no impact on the security of the hash. Two elements in
this definition have an important influence on the security of a hash
function: the choice of the padding rule and the choice of the
\emph{IV}. Figure~\ref{fig:iterated-hash} illustrates the mechanics of a typical
iterated hash function \cite{ref2}.

\begin{figure}[htbp]\centering% Editable reconstruction of Figure 1 in the June 2005 PDF.
\begin{tikzpicture}[>=Latex, every node/.style={font=\small},
 box/.style={draw,minimum width=4cm,minimum height=0.7cm,align=center}]
\node (input) at (0,0.2) {original input: $x$};
\node[box] (pad) at (0,-1) {append padding bits};
\node[box] (length) at (0,-2) {append length block};
\node[draw,fit=(pad)(length),inner sep=0.22cm] (pre) {};
\node[anchor=west] at (2.5,-1.4) {preprocessing};
\node[align=left,anchor=east] at (-2.4,-2.8) {formatted input\\$x=x_1x_2\ldots x_t$};
\node[draw,minimum size=0.9cm] (f) at (0,-4.3) {$f$};
\node[draw,minimum width=3cm,minimum height=2cm] (round) at (0,-4.3) {};
\node[anchor=west,align=left] at (2.5,-4.3) {compression\\function $f$};
\draw (-2.2,-3.25) rectangle (2.2,-5.85);
\node[anchor=west] at (2.5,-5.5) {iterated processing};
\node[box,minimum width=3cm] (g) at (0,-6.7) {$g$};
\node (output) at (0,-7.9) {output: $h(x)=g(H_t)$};
\draw[->] (input)--(pad);
\draw[->] (pad)--(length);
\draw[->] (length)--node[pos=.86,right] {$x_i$}(f);
\draw[->] (f)--node[right,pos=.45] {$H_i$} (0,-5.3);
\draw[->] (0,-5.3)--node[right,pos=.8] {$H_t$}(g);
\draw[->] (0,-5.3)--(-1.8,-5.3)--(-1.8,-4.3)--node[above] {$H_{i-1}$}(f);
\node[anchor=east] at (-2.3,-4.3) {$H_0=IV$};
\draw[->] (g)--(output);
\draw (-5.4,-0.25) rectangle (6,-7.25);
\node[anchor=south east] at (6,-0.25) {hash function $h$};
\end{tikzpicture}
\caption{Iterated Hash Functions}\label{fig:iterated-hash}\end{figure}

It is recommended that the padding rule is unambiguous (i.e., there do
not exist two messages that can be padded to the same padded message);
at the end one should append the length of the message; and the
\emph{IV} should be defined as part of the description of the hash
functions (this is called MD-strengthening, after Merkle and Damgård).
In some cases one can deviate from this rule, but this will make the
hash function less secure and may lead to trivial collisions or second
preimages. Note that the function \emph{f}, the \emph{IV}, and the
padding scheme are publicly known.

Research on hash functions has been focused on the question: which
properties should be imposed on \emph{f} to guarantee that \emph{h}
satisfies certain properties? Two partial answers have been found to
this question. We state them here for the case that \emph{g} is the
identity function. The first result is by Lai and Massey
\cite{ref29} and gives necessary and sufficient conditions for
\emph{f} in order to obtain an ``ideally secure'' hash function
\emph{h}.

\textbf{Theorem 1 (Lai-Massey)} \emph{Assume that the padding contains
the length of the input string, and that the message X (without padding)
contains at least 2 blocks. Then finding a second-preimage for h with a
fixed IV requires 2}\textsuperscript{\emph{n~}} \emph{operations if and
only if finding a second-preimage for f with arbitrarily chosen
H}\textsubscript{\emph{i-1}} \emph{requires
2}\textsuperscript{\emph{n~}}\emph{operations.}

The fact that the condition is necessary is based on the following
argument: if one can find a second-preimage for \emph{f} in
\emph{2}\textsuperscript{\emph{S}} operations (with \emph{s \textless{}
n}), one can find a second-preimage for \emph{h} in \(2^{1+(n+s)/2}\)
operations with a meet-in-the-middle attack \cite{ref5}.

A second result by Damgård \cite{ref15} and independently by
Merkle \cite{ref24} states that for \emph{h} to be a CRHF it
is sufficient that \emph{f} is a collision resistant function.

\textbf{Theorem 2 (Damgård-Merkle)} \emph{Let f be a collision resistant
function mapping l to n bits (with l -- n \textgreater{} 1). If an
unambiguous padding rule is used, the following construction yields a
CRHF:}

\begin{equation}\label{eq:2.2}
\begin{aligned}H_1&=f(0^{n+1}\mathbin\Vert X_1)\\H_i&=f(H_{i-1}\mathbin\Vert1\mathbin\Vert X_i)\quad i=2,3,\ldots,t\\h(X)&=H_{t+1}\end{aligned}
\end{equation}

Where \textbar\textbar{} denotes concatenation of bit strings and
\(0^{n+1}\) indicates a sequence of \emph{n+1} bits, all being zero.

The proof that the resulting function \emph{h} is collision resistant
follows by a simple argument that a collision for \emph{h} would imply a
collision for \emph{f} for some stage \emph{i}. The inclusion of the
length-block, which effectively encodes all messages such that no
encoded input is the tail end of any other encoded input, is necessary
for this reasoning \cite{ref2}.

Recent findings \cite{ref11} show that for a typical
\emph{n}-bit iterated hash function, one can find a second preimage for
a \(2^k\) -block message in about \(k\cdot2^{n/2+1}+2^{n-k+1}\)
operations for long messages, assuming unlimited memory.
